% !TeX root = ../Book4.tex
% 线性约化：在 C={p_2=0} 中沿根空间方向取商。
\begin{figure}[htbp]
	\centering
	\begin{tikzpicture}[x={(1cm,0cm)},y={(0.6cm,0.35cm)},z={(0cm,1cm)},>=Stealth,font=\small]
		\fill[GeometryBlueLight] (-1.7,-1.3,0)--(2.2,-1.3,0)--(2.2,1.7,0)--(-1.7,1.7,0)--cycle;
		\draw[GeometryGray] (-1.7,-1.3,0)--(2.2,-1.3,0)--(2.2,1.7,0)--(-1.7,1.7,0)--cycle;
		\draw[->] (-1.8,0,0)--(2.7,0,0) node[right] {\(q^1\)};
		\draw[->] (0,-1.5,0)--(0,2.1,0) node[above left] {\(p_1\)};
		\draw[->] (0,0,-0.5)--(0,0,2.8) node[above] {\(q^2\)};
		\draw[GeometryBlue,thick] (1,0.6,-0.3)--(1,0.6,2.6);
		\fill[GeometryBlue] (1,0.6,0) circle (2pt);
		\fill[GeometryBlue] (1,0.6,1.2) circle (2pt);
		\fill[GeometryBlue] (1,0.6,2.1) circle (2pt);
		\draw[GeometryBlue,->,thick] (1.24,0.6,2.1)--(1.24,0.6,0.12);
		\node[right] at (1.35,0.6,1.6) {同一个商类};
		\node[below] at (0.8,-1.2,0) {\((q^1,p_1)\) 平面};
		\node[left] at (-0.6,0,2.2) {\(C=\{p_2=0\}\)};
	\end{tikzpicture}
	\caption{三维约束空间 \(C\) 中，平行于根空间 \(\mathbb R\partial_{q^2}\) 的每条直线成为约化空间中的一个点。所画平面 \(q^2=0\) 选取了每个商类的一个代表。}
	\label{g4:fig:0701:linear-reduction}
\end{figure}
